% Generated draft: requires mathematical and source review.
\documentclass[11pt]{article}
\usepackage[margin=25mm]{geometry}
\usepackage{fontspec}
\setmainfont{Latin Modern Roman}
\usepackage{amsmath,amssymb,mathtools,unicode-math}
\setmathfont{Latin Modern Math}
\usepackage{enumitem,needspace,xurl,hyperref}
\setlength{\parindent}{0pt}
\setlength{\parskip}{5pt}
\setlength{\emergencystretch}{3em}
\newcommand{\N}{\mathbb{N}}
\newcommand{\Z}{\mathbb{Z}}
\newcommand{\Q}{\mathbb{Q}}
\newcommand{\R}{\mathbb{R}}
\newcommand{\C}{\mathbb{C}}
\newcommand{\F}{\mathbb{F}}
\newcommand{\A}{\mathbb{A}}
\newcommand{\PP}{\mathbb{P}}
\newcommand{\E}{\mathbb{E}}
\newcommand{\eps}{\varepsilon}
\newcommand{\dd}{\,\mathrm{d}}
\begin{document}
\section*{UnsolvedMath ID: 30002371}
\textbf{Source problem code:} catalogue entry
\subsection{Definitions and mathematical statement}
Let \(n\geq 1\), and let \(M^{2n}\) be a polyhedral Kähler manifold in the sense intended by the supplied source. Let \(S_{\mathrm{met}}\subset M\) be its metric singular locus, namely the points without a neighborhood isometric to a Euclidean open set, and put \(U=M\setminus S_{\mathrm{met}}\). Denote the given flat complex structure on \(U\) by \(J\); it is parallel for the flat metric and makes \(U\) a complex manifold of complex dimension \(n\). The supplied record does not give the full definition of a polyhedral Kähler manifold, including any additional conditions imposed along singular strata. These hypotheses must therefore be understood in the source's sense and checked before treating this as a complete standalone formulation.

\textbf{Complex extension conjecture.} For every such \(M^{2n}\), does there exist a normal complex analytic variety \(X\) of pure complex dimension \(n\) and a homeomorphism \(h:M\to X\) such that \(h(U)\subset X_{\mathrm{reg}}\) is open and \(h|_U:(U,J)\to h(U)\) is biholomorphic, with the further source requirement that \(X\) have ``complex singularities in codimension \(3\)''? Here \(X_{\mathrm{reg}}\) is the complex manifold locus and \(\operatorname{Sing}(X)=X\setminus X_{\mathrm{reg}}\) is the complex analytic singular locus. Normality means that every local analytic ring \(\mathcal O_{X,x}\) is an integrally closed domain. The metric singular locus and the complex analytic singular locus are distinct notions; the extension requirement implies only \(\operatorname{Sing}(X)\subset h(S_{\mathrm{met}})\).

\textbf{Unresolved codimension qualification.} If the quoted phrase means complex codimension at least \(3\), the additional requirement is
\[
\operatorname{codim}_{\C}(Z,X)\geq 3
\quad\text{for every irreducible component }Z\subset\operatorname{Sing}(X).
\]
Equivalently, every such component has complex dimension at most \(n-3\); an empty singular locus satisfies this condition. In particular, this reading requires \(X\) to be nonsingular when \(n\leq 2\). The supplied original statement does not explicitly say ``at least,'' while the normalized statement says that the singularities ``have codimension \(3\).'' An exact-codimension requirement, including whether it applies to every component and whether an empty singular locus is allowed, is therefore not settled by the record. The displayed inequality is a conditional clarification, not a silently imposed correction.
\subsection{Short English statement}
The conjecture asks whether the flat complex structure on the metrically regular part of every polyhedral Kähler manifold extends to a normal complex analytic structure on the entire underlying space, with complex singularities restricted by the source's codimension-\(3\) condition. The precise polyhedral Kähler hypotheses and the meaning of that codimension condition require source review. The supplied dated curation labels the problem partially solved.
\subsection{Sources}
\begin{itemize}
\item Local UnsolvedMath registry: supplied complete record for permanent ID \(30002371\), titled \emph{Complex Extension Conjecture for Polyhedral Kähler Manifolds}, source problem code \texttt{catalogue entry}, in the Oberwolfach Reports Open Problems collection. No individual proposer is identified; the proposed year is recorded as \(2013\).
\item Supplied primary report citation: \emph{Differentialgeometrie im Großen (2014)}, Oberwolfach Reports, pp.\ 1929--1974, DOI \texttt{10.4171/owr/2013/33}. \url{https://doi.org/10.4171/owr/2013/33}. The record describes Oberwolfach Reports as documenting mathematical workshops and their open-problem discussions. Its displayed report year \(2014\) differs from the year \(2013\) in the DOI and proposed-year metadata; this discrepancy is unresolved here.
\item Supplied literature reference: Panov, \emph{polyhedral Kähler manifolds}. \url{https://doi.org/10.2140/gt.2009.13.2205}.
\item Supplied literature reference: de Borbon and Panov, \emph{polyhedral Kähler metrics}. \url{https://arxiv.org/abs/2510.17447}.
\item Supplied literature reference: \emph{Oberwolfach report on polyhedral Kähler manifolds}. \url{https://oa.tib.eu/renate/server/api/core/bitstreams/2771cf96-e925-492b-8b7c-0481c8c56da5/content}.
\item Literature assessment dated \texttt{2026-08-21}: the supplied curation labels the problem \emph{partially solved} and reports that Panov proves the extension and complex-surface structure in real dimension four, with further special constructions in higher dimensions, while the all-dimensional complex conjecture remains open. This is supplied dated curation derived from the OWR record, not an independently verified assessment; the time-sensitive status and cited claims require independent verification.
\item The registry labels the statement \texttt{corrected\_verified} and reports that notation, typography, and wording were normalized without changing mathematical content. This is supplied editorial metadata, not independent mathematical verification of this draft.
\end{itemize}
\end{document}
